% header.tex — Pandoc PDF setup (Farben, Boxen, Code-Umbruch, Seitenränder)

\usepackage[a4paper,top=2cm,bottom=2cm,left=2cm,right=2cm]{geometry}

% Bessere Darstellung von Quotes/Backticks in verbatim
\usepackage{upquote}

% Bessere Code-/Verbatim-Umbrüche
\usepackage{fvextra}
\DefineVerbatimEnvironment{Highlighting}{Verbatim}{
  breaklines,
  breakanywhere,
  fontsize=\small,
  commandchars=\\\{\}
}

\usepackage{xcolor}
\usepackage{hyperref}
\usepackage{titlesec}
\usepackage[most]{tcolorbox}

\AtBeginDocument{%
  \hypersetup{
    colorlinks=true,
    allcolors=blue,
    pdfborderstyle={/S/U/W 1},
    pdfborder={0 0 1}
  }%
}

% --------------------------------------------------
% Farben
% --------------------------------------------------

% Überschriften
\definecolor{SectionBlue}{HTML}{1F4E79}
\definecolor{SubsectionTeal}{HTML}{0F766E}
\definecolor{SubsubsectionSlate}{HTML}{334155}

% Boxen
\definecolor{InfoBack}{HTML}{EEF6FF}
\definecolor{InfoFrame}{HTML}{2563EB}

\definecolor{NoteBack}{HTML}{F8FAFC}
\definecolor{NoteFrame}{HTML}{64748B}

\definecolor{ImportantBack}{HTML}{FFF7D6}
\definecolor{ImportantFrame}{HTML}{D97706}

\definecolor{WarningBack}{HTML}{FEECEC}
\definecolor{WarningFrame}{HTML}{DC2626}

\definecolor{TipBack}{HTML}{ECFDF5}
\definecolor{TipFrame}{HTML}{059669}

\definecolor{ExampleBack}{HTML}{F5F3FF}
\definecolor{ExampleFrame}{HTML}{7C3AED}

\definecolor{StepBack}{HTML}{F3F4F6}
\definecolor{StepFrame}{HTML}{4B5563}

\definecolor{DefinitionBack}{HTML}{EFFCF6}
\definecolor{DefinitionFrame}{HTML}{0F766E}

\definecolor{SummaryBack}{HTML}{FFF1F2}
\definecolor{SummaryFrame}{HTML}{E11D48}

\definecolor{ChecklistBack}{HTML}{F0FDF4}
\definecolor{ChecklistFrame}{HTML}{16A34A}

\definecolor{ReferenceBack}{HTML}{FAF5FF}
\definecolor{ReferenceFrame}{HTML}{9333EA}

% --------------------------------------------------
% Überschriften
% --------------------------------------------------

\titleformat{\section}
  {\Large\bfseries\color{SectionBlue}}
  {\thesection}
  {0.75em}
  {}

\titleformat{\subsection}
  {\large\bfseries\color{SubsectionTeal}}
  {\thesubsection}
  {0.75em}
  {}

\titleformat{\subsubsection}
  {\normalsize\bfseries\color{SubsubsectionSlate}}
  {\thesubsubsection}
  {0.75em}
  {}

% --------------------------------------------------
% Allgemeiner Stil für Boxen
% --------------------------------------------------

\tcbset{
  docbox/.style={
    enhanced,
    breakable,
    boxrule=0.8pt,
    arc=2mm,
    left=2mm,
    right=2mm,
    top=1mm,
    bottom=1mm,
    fonttitle=\bfseries
  }
}

% --------------------------------------------------
% Universelle Boxen
% --------------------------------------------------

\newtcolorbox{infobox}{
  docbox,
  colback=InfoBack,
  colframe=InfoFrame,
  title=Info
}

\newtcolorbox{notebox}{
  docbox,
  colback=NoteBack,
  colframe=NoteFrame,
  title=Hinweis
}

\newtcolorbox{importantbox}{
  docbox,
  colback=ImportantBack,
  colframe=ImportantFrame,
  title=Wichtig
}

\newtcolorbox{warningbox}{
  docbox,
  colback=WarningBack,
  colframe=WarningFrame,
  title=Achtung
}

\newtcolorbox{tipbox}{
  docbox,
  colback=TipBack,
  colframe=TipFrame,
  title=Tipp
}

\newtcolorbox{examplebox}{
  docbox,
  colback=ExampleBack,
  colframe=ExampleFrame,
  title=Beispiel
}

\newtcolorbox{stepbox}{
  docbox,
  colback=StepBack,
  colframe=StepFrame,
  title=Vorgehen
}

\newtcolorbox{definitionbox}{
  docbox,
  colback=DefinitionBack,
  colframe=DefinitionFrame,
  title=Definition
}

\newtcolorbox{summarybox}{
  docbox,
  colback=SummaryBack,
  colframe=SummaryFrame,
  title=Zusammenfassung
}

\newtcolorbox{checklistbox}{
  docbox,
  colback=ChecklistBack,
  colframe=ChecklistFrame,
  title=Checkliste
}

\newtcolorbox{referencebox}{
  docbox,
  colback=ReferenceBack,
  colframe=ReferenceFrame,
  title=Verweis
}

% --------------------------------------------------
% Optional: jede Section auf neuer Seite
% Auskommentieren, wenn du Fließlayout willst
% --------------------------------------------------
% \newcommand{\sectionbreak}{\clearpage}


% --------------------------------------------------
% Mitschrift aus Unterricht
% --------------------------------------------------

\definecolor{UnterrichtBack}{HTML}{F0FDF4}
\definecolor{UnterrichtFrame}{HTML}{16A34A}

\newtcolorbox{unterrichtbox}{
  docbox,
  colback=UnterrichtBack,
  colframe=UnterrichtFrame,
  title=Mitschrift aus Unterricht
}


% --------------------------------------------------
% Bilder nicht frei schweben lassen
% --------------------------------------------------

\usepackage{float}

\makeatletter
\def\fps@figure{H}
\makeatother